% Diagramma TikZ per regole di bonifica ed evoluzione dello schema (essenziale e pulito)
\begin{tikzpicture}[
    font=\sffamily,
    >=Stealth,
    node distance=0.8cm,
    ruleBox/.style={
        draw=brandBorderStrong,
        fill=white,
        rounded corners=4pt,
        line width=0.8pt,
        minimum width=6.2cm,
        minimum height=1.05cm,
        align=left,
        inner xsep=8pt,
        inner ysep=4pt
    },
    schemaBox/.style={
        draw=brandBorderStrong,
        fill=white,
        rounded corners=4pt,
        line width=0.9pt,
        minimum width=5.6cm,
        minimum height=1.35cm,
        align=center,
        inner sep=6pt
    },
    flowArrow/.style={
        ->,
        line width=1.4pt,
        color=brandPrimary
    }
]

    % ==================== COLONNA SINISTRA: 4 REGOLE DI BONIFICA ====================
    \node[font=\footnotesize\bfseries\color{brandPrimary}, anchor=west] at (-6.8, 3.4) {4 Regole Fisiologiche di Bonifica};

    % Regola 1: Outlier
    \node[ruleBox, draw=brandCoral, fill=brandCoral!5] (r1) at (-3.7, 2.5) {
        \textbf{\color{brandCoralInk}1. Mascheramento Outlier}\\
        {\scriptsize $x \notin \text{TRACK\_SPECS} \implies \texttt{*\_outlier = True} \land x = \text{NaN}$}
    };

    % Regola 2: Interpolazione
    \node[ruleBox, draw=brandPrimary, fill=brandPrimary!5, below=0.35cm of r1] (r2) {
        \textbf{\color{brandPrimary}2. Interpolazione Lineare Gap Brevi}\\
        {\scriptsize $\Delta t \le 5\text{ s} \implies \text{interpolazione}$ (\texttt{HR\_imputed = True})}
    };

    % Regola 3: LOCF
    \node[ruleBox, draw=brandPeachInk, fill=brandPeach!8, below=0.35cm of r2] (r3) {
        \textbf{\color{brandPeachInk}3. Propagazione NIBP (LOCF)}\\
        {\scriptsize $\Delta t \le 300\text{ s} \implies \text{mantieni costante}$ (\texttt{SBP\_imputed = True})}
    };

    % Regola 4: Decadimento
    \node[ruleBox, draw=brandTeal, fill=brandTeal!5, below=0.35cm of r3] (r4) {
        \textbf{\color{brandTeal}4. Contatore di Freschezza Misura}\\
        {\scriptsize $\texttt{nibp\_age\_s} = t - t_{\text{last\_cuff}} \quad (\text{sentinella } 9999.0)$}
    };

    % ==================== COLONNA DESTRA: EVOLUZIONE DELLO SCHEMA ====================
    \node[font=\footnotesize\bfseries\color{brandDark}, anchor=west] at (1.0, 3.4) {Evoluzione Formale dello Schema};

    % Schema 1: Bronze
    \node[schemaBox, draw=brandCoral, fill=brandCoral!6] (s1) at (3.8, 2.3) {
        {\textbf{\color{brandCoralInk}Bronze Parquet (6 colonne)}}\\[2pt]
        {\scriptsize \texttt{Time}, 5 tracce biometriche grezze}
    };

    % Schema 2: Silver
    \node[schemaBox, draw=brandPrimary, fill=brandPrimary!6, below=0.9cm of s1] (s2) {
        {\textbf{\color{brandPrimary}Silver Parquet (18 colonne)}}\\[2pt]
        {\scriptsize 6 vitali + 5 outlier + 5 imputed + \texttt{age} + \texttt{dep}}
    };

    % Schema 3: Gold
    \node[schemaBox, draw=brandTeal, fill=brandTeal!6, below=0.9cm of s2] (s3) {
        {\textbf{\color{brandTeal}Gold Documento BSON (4 sezioni)}}\\[2pt]
        {\scriptsize \texttt{timestamp}, \texttt{metadata}, \texttt{metrics}, \texttt{flags} (NaN omessi)}
    };

    % Frecce verticali di transizione dello schema
    \draw[flowArrow] (s1) -- (s2)
        node[midway, right=3pt, font=\scriptsize\bfseries\color{brandPrimary}] {+12 flag \& feature};

    \draw[flowArrow, color=brandTeal] (s2) -- (s3)
        node[midway, right=3pt, font=\scriptsize\bfseries\color{brandTeal}] {BSON (-40 pt null)};

\end{tikzpicture}
